\section{\texorpdfstring{\code{volatile}}{volatile}}
\label{sec:volatile}

So far the code has only touched ordinary memory. This file touches \textbf{hardware}, and at that
point one of the compiler's fundamental assumptions stops holding.

The compiler assumes that nobody but the program changes a variable. On that assumption the
following loop is infinite, and the optimiser is allowed to remove it:

\begin{cppcode}
while ((SPI0.INTFLAGS & SPI_IF_bm) == 0U)
{
}
\end{cppcode}

\code{INTFLAGS} is read once, the result is put in a register, and since nothing in the loop
changes it the condition is constant. With \code{-O2} the result is either a loop that never ends
or a loop that is optimised away entirely.

But \code{INTFLAGS} \emph{does change} --- by the peripheral, without the program doing anything.
\code{volatile} is how you say that: \emph{re-read this variable every time; do not merge reads, do
not move them, do not remove them.}

In \code{<avr/io.h>} the peripheral registers are already declared \code{volatile}. You do not have
to write the qualifier --- but you do have to know that it is there, and why.

\subsection{What \texorpdfstring{\code{volatile}}{volatile} is not}

It is \textbf{not} a synchronisation mechanism. It guarantees neither atomicity nor any ordering
against anything but other \code{volatile} accesses. In threaded code on a host computer
\code{volatile} is almost always the wrong tool; \code{std::atomic} is the right one.

Here it is about something else: a register that a peripheral changes under the program's nose.
That is precisely the case \code{volatile} exists for.

\subsection{Where it may appear}

\textbf{Only in the transport's implementation file.}

That is the whole point of the seam in \kapref{ch:byte}. Everything above it --- \code{Spi},
\code{Stub}, \code{Interface}, \code{EchoNode} --- is ordinary, portable C++ that compiles and runs
on your laptop. A reviewer who sees \code{volatile} in any other file should ask why.

\section{The SPI peripheral}
\label{sec:periph}

Four registers are enough for a polled master: one for enable and clock settings, one for the
mode, one for status flags, and one for data.

\begin{cppcode}
void AvrSpiTransport::init() noexcept
{
    // Route SPI0 to PORTC (ALT1): PC0 = MOSI, PC1 = MISO, PC2 = SCK, PC3 = SS.
    PORTMUX.SPIROUTEA = PORTMUX_SPI0_ALT1_gc;

    // MOSI, SCK and SS are outputs; MISO is an input. SS idles high.
    PORTC.DIRSET = PIN0_bm | PIN2_bm | PIN3_bm;
    PORTC.DIRCLR = PIN1_bm;
    PORTC.OUTSET = PIN3_bm;

    // Master, MSB first, 4 MHz / 4 = 1 MHz, enabled.
    SPI0.CTRLA = SPI_MASTER_bm | SPI_PRESC_DIV4_gc | SPI_ENABLE_bm;

    // SPI mode 0, and software-controlled slave select.
    SPI0.CTRLB = SPI_MODE_0_gc | SPI_SSD_bm;
}
\end{cppcode}

Five lines, and three of them are traps:

\textbf{The pin mux.} Without that line the peripheral sits on its default mux, and the four wires
you connected are completely silent. The symptom is the most confusing there is: the program runs,
no errors are reported, and the logic analyser shows nothing at all. \textbf{Set it first.}

\textbf{Slave select disable.} Without it the peripheral monitors the \code{SS} pin, and a low
level there makes the master believe that another master has taken over the bus: it leaves master
mode. But here it is \emph{your own code} that pulls \code{SS} low on every \code{select()}.
Without the bit, the driver shoots itself in the foot on the first transaction.

\textbf{The prescaler.} The protocol guarantees nothing above 1~MHz. Do not run faster just
because you can.



\subsection{The transport}

\begin{cppcode}
void AvrSpiTransport::select() noexcept   { PORTC.OUTCLR = PIN3_bm; }
void AvrSpiTransport::deselect() noexcept { PORTC.OUTSET = PIN3_bm; }

std::uint8_t AvrSpiTransport::transfer(const std::uint8_t byte) noexcept
{
    SPI0.DATA = byte;

    while ((SPI0.INTFLAGS & SPI_IF_bm) == 0U)
    {
    }

    return SPI0.DATA;
}
\end{cppcode}

\textbf{Reading \code{SPI0.DATA} clears the flag.} Skip it and the next \code{transfer()} sees a
flag that is still set, returns immediately, and reads a byte that has not had time to arrive. The
symptom is data lagging \textbf{one byte behind} through the whole transaction, which is hard to
see and easy to mistake for a bit-order error.

\textbf{No delay is needed around \code{SS}.} The protocol requires 60~ns of setup and hold; a
single port write on a 4~MHz MCU is already several times slower than that.

\section{MVIO: two voltages in the same chip}
\label{sec:mvio}

This is the most elegant detail in the whole design.

\textbf{The problem:} the MCU board runs at 5~V. The FPGA board's GPIO is 3.3~V and does not
tolerate 5~V. The usual solution is a level shifter on every wire --- one more component, four more
connections, and four more things that can be in the wrong place.

\textbf{The solution:} the AVR-DB family has \textbf{MVIO} (\emph{Multi-Voltage I/O}).
\code{PORTC} is not fed from \code{VDD} like the rest of the chip, but from a pin of its own:
\code{VDDIO2}.

Connect \code{VDDIO2} to the FPGA board's 3.3~V rail, and the whole of \code{PORTC} becomes a
3.3-volt domain --- while the core and the other ports run unchanged at 5~V. \textbf{No level
shifter.}

That is why the pin configuration is part of the contract and not each group's own choice: switch
to the default mux and the wires end up in the \code{VDD} domain, and then the level shifter is
needed again --- or the FPGA input breaks.

\begin{aside}
\note \code{VDDIO2} has to be powered \textbf{before} \code{PORTC} does anything at all. If it is
not, the port's behaviour is undefined. The MVIO status register says whether it is, and checking
it at start-up is the single most valuable line of code ahead of bring-up: a forgotten
\code{VDDIO2} connection looks exactly like a broken SPI implementation, and that check tells the
two cases apart in a second.
\end{aside}

\section{Freestanding C++}
\label{sec:freestanding}

The firmware is built for a \textbf{freestanding} target. Part of C++ does not come along:

\begin{booktable}{@{}lL@{}}
\thead{Not available} & \thead{What you do instead} \\ \midrule
The heap (\code{new}, \code{std::vector}) & Statically allocated objects; everything has a known
  size. \\
Exceptions & Return codes --- which the whole interface is already built on. \\
RTTI (\code{dynamic_cast}) & Virtual functions are enough; you never cast downwards. \\
\code{<iostream>} & UART output, or an LED. \\
\end{booktable}

\textbf{Note that none of that affects the design.} The interface returns \code{bool} instead of
throwing. \code{Frame} is an aggregate with a static size. No class uses the heap. That was not
luck --- the project's coding conventions are written for a freestanding target from the start, and
this is the chapter where that pays off.
